\documentclass[10pt,a4paper]{article}
\usepackage[margin=2.0cm]{geometry}
\usepackage{fontspec}
\usepackage{amsmath,amssymb,amsfonts,amsthm,mathtools}
\usepackage{booktabs,tabularx,colortbl,array,multirow}
\usepackage[table]{xcolor}
\usepackage{fancyhdr}
\usepackage{titlesec}
\usepackage{microtype}
\usepackage{hyperref}

\hypersetup{
    colorlinks=true,
    linkcolor=blue!70!black,
    citecolor=blue!70!black,
    urlcolor=blue!70!black
}

\pagestyle{fancy}
\fancyhf{}
\rhead{\footnotesize \textbf{NEXVIA AEGIS V9-M4 Final Research Audit}}
\lhead{\footnotesize Independent Statistical \& Architecture Review}
\rfoot{\footnotesize Page \thepage}
\lfoot{\footnotesize Document ID: AEGIS-V9-M4-AUD-FINAL}

\titleformat{\section}{\large\bfseries}{\thesection}{1em}{}
\titleformat{\subsection}{\normalsize\bfseries}{\thesubsection}{1em}{}

\begin{document}

\begin{center}
    {\LARGE \textbf{NEXVIA AEGIS V9-M4 Independent Statistical Audit}}\\[6pt]
    {\Large \textbf{Comprehensive Scientific Audit Report: Confirmatory Holdout Evaluation Harness, Statistical Formulations, and Pre-Execution Readiness}}\\[10pt]
    \textbf{Autonomous Verification Subsystem (Google DeepMind Antigravity)}\\
    \textit{Git Branch: v9.0-runtime-assurance \quad $\cdot$ \quad Commit: 083ade9 \quad $\cdot$ \quad Date: October 25, 2026}\\[12pt]
\end{center}

\noindent\fbox{
\parbox{0.97\textwidth}{
\textbf{NON-AUTHORIZATION \& GOVERNANCE STATEMENT:}\\
This document reports an independent statistical, mathematical, and implementation audit of the NEXVIA AEGIS V9-M4 confirmatory holdout evaluation harness. \textbf{No real holdout execution was authorized or performed as part of this audit.} All empirical tests reported herein are based strictly on synthetic mock records or unit regression baselines. The reserved holdout seed blocks ($73000..74439$ for mixture holdout and $75000..75119$ for scenario holdout) and their corresponding simulation paths remain completely uninstantiated and 100\% untouched. This audit report does NOT constitute authorization to execute V9-M4.
}
}

\vspace{10pt}

\begin{abstract}
\noindent This comprehensive scientific audit report presents the complete findings, mathematical derivations, numerical validations, and architectural reviews conducted on the NEXVIA AEGIS V9-M4 confirmatory holdout evaluation harness (\texttt{experiments/v9\_evaluate\_holdout.py}). We audit the statistical foundations of preregistered Hypotheses $H_1$ through $H_5$, evaluating exact Clopper--Pearson confidence bound inversion, Berger Intersection-Union Tests (IUT), step-down Holm--Šidák familywise error control, exact McNemar testing, Newcombe paired difference intervals (Method 10), and empirical coverage guardrails. Our audit confirms that the evaluation runner achieves high mathematical fidelity in its core estimators, passes all 125 full repository regression tests, enforces robust single-shot overwrite guards, and maintains complete isolation of reserved holdout partitions. However, our investigation identified one high-severity statistical discrepancy (\textbf{Finding M4-F02}): in evaluating the multi-scenario safety envelope ($H_5$), the harness calculates binomial $p$-values on accepted decisions ($M_s$) with a punitive fallback of $p=1.0$ under total abstention, rather than on unconditional incidence ($N_s = 120$) as specified in the preregistered protocol. This defect inadvertently penalizes safe runtime abstention in catastrophic fault scenarios, guaranteeing that the Berger IUT would fail even under zero false successes ($FS=0$). We provide complete mathematical proofs, synthetic boundary data, and an actionable remediation specification. We recommend a status of \textbf{\texttt{M4\_STATISTICAL\_AUDIT\_PASS\_WITH\_QUALIFICATIONS}}, requiring the remediation of Finding M4-F02 and Finding M4-F03 prior to formal production holdout authorization.
\end{abstract}

\vspace{6pt}
\noindent \textbf{Keywords:} Certified Runtime Assurance, Selective Risk Control, Clopper--Pearson Exact Bounds, Berger Intersection-Union Test, Step-Down Holm--Šidák, Paired McNemar Testing, Holdout Data Governance.

\section{Introduction \& Audit Objectives}
Milestone V9-M4 represents the confirmatory holdout evaluation of the calibrated runtime assurance policy:
\begin{equation}
\hat{\pi} = \text{V9\_CAND\_1\_TIER1\_STRICTEST}
\end{equation}
Under research governance protocols, holdout evaluation is strictly unrepeatable. To guarantee that the confirmatory evaluation is methodologically sound, reproducible, and immune to post-hoc adjustments, this independent pre-execution audit evaluates:
\begin{enumerate}
    \item Traceability between preregistered hypotheses ($H_1$–$H_5$) and executable code.
    \item Exactness and boundary behavior of statistical estimators and confidence bounds.
    \item Behavior under edge cases (e.g.\ $M=0$, small $M$, extreme fault regimes).
    \item Robustness of execution guards, single-shot protections, and seed isolation.
    \item Algorithmic regressions across the full repository test suite.
\end{enumerate}

\section{Traceability Matrix: Hypotheses to Implementation}
Table~\ref{tab:traceability} summarizes the complete hypothesis traceability across the harness.

\begin{table}[h!]
\centering
\scriptsize
\begin{tabularx}{\textwidth}{p{0.6cm}p{1.8cm}p{1.3cm}Xp{2.2cm}p{2.3cm}p{2.1cm}c}
\toprule
\textbf{Hyp.} & \textbf{Description} & \textbf{Partition} & \textbf{Estimand \& Definition} & \textbf{Statistical Test} & \textbf{Pass Threshold} & \textbf{Implementation} & \textbf{Status} \\
\midrule
$H_1$ & \textbf{Selective Risk Certification} & $\mathcal{D}_{\mathrm{test}}^{\mathrm{mix}}$ ($N=1440$) & $R_{\mathrm{sel}}(\hat{\pi}) = \mathbb{P}(FS=1 \mid A=1)$ on $\mathcal{P}_{12}$ & 1-sided exact Clopper--Pearson 95\% UCB & $M \ge 59 \land U_{\mathrm{CP},1} \le 0.05 \land \hat{C} \ge 0.35$ & \texttt{evaluate\_mixture\_holdout} (lines 539--566) & \textbf{VERIFIED} \\
\addlinespace
$H_2$ & \textbf{Unconditional FS Reduction} & $\mathcal{D}_{\mathrm{test}}^{\mathrm{mix}}$ ($N=1440$) & $\Delta R_{\mathrm{inc}} = R_{\mathrm{inc}}(\hat{\pi}) - R_{\mathrm{inc}}(\pi_{\mathrm{base}})$ & Exact 2-sided McNemar + Newcombe paired CI & $p_{\mathrm{McNemar}} < 0.05 \land \hat{R}_{\mathrm{inc}} < \hat{R}_{\mathrm{inc, base}}$ & \texttt{evaluate\_mixture\_holdout} (lines 568--615) & \textbf{VERIFIED} \\
\addlinespace
$H_3$ & \textbf{TSR Noninferiority} & $\mathcal{D}_{\mathrm{test}}^{\mathrm{scen}}$ ($N=1200$) & $\Delta \mathrm{TSR} = \mathrm{TSR}(\hat{\pi}) - \mathrm{TSR}(\pi_{\mathrm{V7}})$ on $\mathcal{S}_{10}$ & 1-sided 95\% paired difference Student's $t$ LCB & $\mathrm{LCB}_{95}(\Delta\mathrm{TSR}) > -0.05 \land \widehat{\mathrm{TSR}}_{\mathrm{NOMINAL}} \ge 0.94$ & \texttt{evaluate\_scenario\_holdout} (lines 740--767) & \textbf{VERIFIED} \\
\addlinespace
$H_4$ & \textbf{Coverage Floor Precondition} & $\mathcal{D}_{\mathrm{test}}^{\mathrm{mix}}$ ($N=1440$) & $C(\hat{\pi}) = \mathbb{P}(A=1)$ on $\mathcal{P}_{12}$ & Point estimate $\hat{C} = M/N$ + Wilson score CI & $\hat{C}(\hat{\pi}; \mathcal{P}_{12}) \ge 0.35$ (target $\ge 0.45$) & \texttt{evaluate\_mixture\_holdout} (lines 545--555) & \textbf{VERIFIED} \\
\addlinespace
$H_5$ & \textbf{Worst-Case Scenario Safety} & $\mathcal{D}_{\mathrm{test}}^{\mathrm{scen}}$ ($12 \times 120$) & $R_{\mathrm{inc}, s} = \mathbb{P}(FS=1 \mid \text{scenario } s)$ for $s \in \mathcal{S}_{12}$ & Berger Intersection-Union Test + Holm--Šidák bounds & $\max_{s \in \mathcal{S}_{12}} p_{\mathrm{inc}, s} \le 0.05$ & \texttt{evaluate\_scenario\_holdout} (lines 652--774) & \textbf{FINDING M4-F02} \\
\bottomrule
\end{tabularx}
\caption{Traceability matrix linking preregistered protocol hypotheses to evaluation harness code.}
\label{tab:traceability}
\end{table}

\section{Summary of Technical Findings by Hypothesis}

\subsection{Hypothesis $H_1$: Selective False-Success Risk Control}
\begin{itemize}
    \item \textbf{Mathematical Derivation of Sample Size Floor:} We derived that for zero false successes ($K=0$), $U_{\mathrm{CP}, 1}(0, M; 0.05) = 1 - 0.05^{1/M}$. Setting $1 - 0.05^{1/M} \le 0.05$ yields $M \ge \ln(0.05)/\ln(0.95) \approx 58.404 \implies M_{\min} = 59$.
    \item \textbf{Anti-Gaming Immunity:} In Case C06 ($M = 59, K = 1$), the empirical point estimate is $1.69\% < 5.0\%$, but the harness correctly rejects certification because the sampling uncertainty expands the bound to $7.79\%$.
    \item \textbf{Clean Edge Case Handling:} For $M=0$, selective risk is properly marked as \texttt{UNDEFINED}, avoiding division-by-zero errors.
\end{itemize}

\subsection{Hypothesis $H_5$: Estimand Inversion in Berger IUT (\textbf{Finding M4-F02})}
The preregistered protocol specifies that the multi-scenario safety envelope is tested via Berger IUT over \textbf{unconditional false-success incidence}:
\begin{equation}
\mathcal{H}_{0,5}: \bigcup_{s=1}^{12} \{R_{\mathrm{inc}, s} > 0.05\} \quad\text{vs.}\quad \mathcal{H}_{1,5}: \bigcap_{s=1}^{12} \{R_{\mathrm{inc}, s} \le 0.05\}
\end{equation}
In \texttt{experiments/v9\_evaluate\_holdout.py}, line 698 computes binomial $p$-values over accepted decisions ($M_s$) with a fallback of $p=1.0$ when $M_s=0$.
\begin{itemize}
    \item In catastrophic fault scenarios (e.g.\ \texttt{FAULTY\_REFERENCE\_SENSOR}), safe controllers abstain on 100\% of episodes ($M_s=0, K_s=0$). Unconditional incidence is $0/120 = 0.0\%$, which achieves exact binomial $p = (0.95)^{120} \approx 0.0021 \le 0.05$.
    \item However, because line 703 assigns $p=1.0$, the Berger IUT $\max p_s \le 0.05$ is guaranteed to fail!
    \item \textbf{Remediation:} Update line 698 to compute $p_{\mathrm{inc}, s} = \mathrm{binom.cdf}(fs_s, n_s, 0.05)$ on $N_s=120$ for the IUT, while preserving descriptive selective risk reporting.
\end{itemize}

\subsection{Hypotheses $H_2$ \& $H_3$: Pairing Alignment (\textbf{Finding M4-F03})}
Exact McNemar testing and Newcombe paired difference intervals are mathematically sound. However, the evaluation loops rely on list ordering via \texttt{zip()} without explicitly verifying \texttt{assert r["seed"] == v["seed"] and r["scenario"] == v["scenario"]}. Adding explicit equality assertions will guarantee robust alignment.

\subsection{Hypothesis $H_4$: Coverage Guardrail}
The evaluation harness correctly enforces the point estimate $\hat{C} \ge 0.35$ as an operational utility floor. At $M=503$, it fails; at $M=504$, it passes. The anti-gaming hierarchy prevents certification if coverage is below 35\%.

\section{Consolidated Finding Severity Register}
Table~\ref{tab:severity} details the formal finding register for Milestone V9-M4.

\begin{table}[h!]
\centering
\small
\begin{tabularx}{\textwidth}{p{1.2cm}cp{3.8cm}Xp{2.2cm}}
\toprule
\textbf{ID} & \textbf{Severity} & \textbf{Title / Area} & \textbf{Finding Description \& Practical Impact} & \textbf{Remediation Requirement} \\
\midrule
\textbf{M4-F01} & Low & Protocol Evolution in Hypothesis Indexing ($H_2$ vs $H_4$) & Early draft narrative labeled $H_2$ as core-10 mixture risk and $H_4$ as compound McNemar. Finalized M4 review formalized $H_2$ as unconditional FS reduction and $H_4$ as coverage floor. Harness correctly follows finalized M4 review. & Reconciled in audit documentation. \\
\addlinespace
\textbf{M4-F02} & \textbf{High} & Estimand Inversion in Berger Intersection-Union Test ($H_5$) & Protocol specifies IUT over unconditional incidence $R_{\mathrm{inc}, s}$ on $N_s=120$. Harness evaluates selective risk on $M_s$, setting $p=1.0$ when $M_s=0$. This penalizes safe abstention in catastrophic fault scenarios and guarantees IUT failure. & Update IUT $p$-value to $p_{\mathrm{inc}, s} = \mathrm{binom.cdf}(fs, n_s, 0.05)$ prior to execution. \\
\addlinespace
\textbf{M4-F03} & Medium & Implicit Key Ordering in Paired Comparisons ($H_2, H_3$) & Paired loops in lines 555--559 and 755--760 use \texttt{zip()} without asserting matching seeds and scenarios. If record arrays are permuted, mismatched episodes are paired silently. & Add explicit assertions: \texttt{assert r["seed"] == v["seed"] and r["scenario"] == v["scenario"]}. \\
\bottomrule
\end{tabularx}
\caption{Consolidated finding severity register for Milestone V9-M4 preflight audit.}
\label{tab:severity}
\end{table}

\section{Regression Testing and Partition Isolation Audit}
\begin{itemize}
    \item \textbf{Preflight Synthetic Suite:} 14 passed in 0.99 seconds (100\% pass rate).
    \item \textbf{Full Repository Regression Suite:} 125 passed in 30.72 seconds (100\% pass rate, 0 failures).
    \item \textbf{Filesystem Cleanliness:} Verified that no real holdout manifests or output directories exist.
    \item \textbf{Fail-Closed Double Lock:} Confirmed that execution without authorization token raises \texttt{RuntimeError}, and live execution is blocked in preflight mode.
\end{itemize}

\section{Audit Conclusions \& Recommendations}
The V9-M4 confirmatory evaluation harness represents an exceptionally rigorous, well-defended implementation of certified runtime assurance. Upon applying the remediation for Finding M4-F02 and Finding M4-F03, the harness will be fully ready for formal authorization review.

\noindent \textbf{Overall Audit Recommendation:} \texttt{M4\_STATISTICAL\_AUDIT\_PASS\_WITH\_QUALIFICATIONS}\\
\noindent \textbf{Holdout Execution Status:} \texttt{M4\_EXECUTION\_STRICTLY\_BLOCKED\_PENDING\_REMEDIATION}

\begin{thebibliography}{9}
\bibitem{clopper1934} C. J. Clopper and E. S. Pearson, ``The use of confidence or fiducial limits illustrated in the case of the binomial,'' \textit{Biometrika}, vol. 26, no. 4, pp. 404--413, 1934.
\bibitem{berger1982} R. L. Berger, ``Multiparameter hypothesis testing and acceptance sampling,'' \textit{Technometrics}, vol. 24, no. 4, pp. 295--300, 1982.
\bibitem{holm1979} S. Holm, ``A simple sequentially rejective multiple test procedure,'' \textit{Scandinavian Journal of Statistics}, vol. 6, no. 2, pp. 65--70, 1979.
\bibitem{shidak1967} Z. Šidák, ``Rectangular confidence regions for the means of multivariate normal distributions,'' \textit{Journal of the American Statistical Association}, vol. 62, no. 318, pp. 626--633, 1967.
\bibitem{mcnemar1947} Q. McNemar, ``Note on the sampling error of the difference between correlated proportions or percentages,'' \textit{Psychometrika}, vol. 12, no. 2, pp. 153--157, 1947.
\bibitem{newcombe1998} R. G. Newcombe, ``Improved confidence intervals for the difference between binomial proportions based on paired data,'' \textit{Statistics in Medicine}, vol. 17, no. 22, pp. 2635--2650, 1998.
\bibitem{wilson1927} E. B. Wilson, ``Probable inference, the law of succession, and statistical inference,'' \textit{Journal of the American Statistical Association}, vol. 22, no. 158, pp. 209--212, 1927.
\bibitem{casella2002} G. Casella and R. L. Berger, \textit{Statistical Inference}, 2nd ed., Duxbury Press, 2002.
\bibitem{ieee730} IEEE Standard for Software Quality Assurance Processes, \textit{IEEE Std 730-2014}, IEEE, 2014.
\end{thebibliography}

\end{document}
